\documentclass[11pt]{article}
\usepackage[margin=1in]{geometry}
\usepackage{amsmath}
\usepackage{array}
\usepackage{booktabs}
\usepackage{longtable}

\title{Memo 02: Reconstructed Delaware River Test Scenarios}
\author{Codex review for FVCOM-MP test-case recovery}
\date{\today}

\begin{document}
\maketitle

\section{Purpose}
This memo reconstructs the intent of the three Delaware River test cases stored in
\texttt{FVCOM\_MP\_test\_run}: cases \texttt{a}, \texttt{b}, and \texttt{c}. The goal is to recover
what is shared across the setups, what actually differs among them, and what those differences imply
for later mass-conservation testing of the floc-enabled plastic model.

\section{Recovered Test Bundle}
\subsection{Shared Domain and Forcing Files}
The common input bundle in \texttt{FVCOM\_MP\_test\_run/INPUT} defines a single Delaware River Estuary
configuration with:
\begin{itemize}
\item 68416 mesh nodes and 30 sigma layers from \texttt{waterPACT\_DRE\_its.nc};
\item 100 open-boundary nodes in \texttt{waterPACT\_obc\_2019.nc} and \texttt{waterPACT\_tsobc\_2019.nc};
\item 9 river entries in \texttt{waterPACT\_riv\_floc\_MP.nml}: 6 Delaware River points
(\texttt{DR\_1}--\texttt{DR\_6}) and 3 Schuylkill River points (\texttt{SR\_1}--\texttt{SR\_3});
\item hourly wind and pressure forcing for calendar year 2019;
\item a river forcing file that spans \texttt{2019/01/01 00:00:01} to \texttt{2020/01/01 00:00:01}.
\end{itemize}

The model runs themselves are set from modified Julian day 58485.0 to 58848.0, i.e.
\texttt{2019-01-02 00:00:00} through \texttt{2019-12-31 00:00:00}. Therefore the shared forcing files
provide roughly a one-day pad before the model start.

\subsection{River Source Characteristics}
The river forcing file \texttt{waterPACT\_riv\_floc\_MP.nc} contains:
\begin{itemize}
\item \texttt{river\_flux} in \(\mathrm{m^3\,s^{-1}}\);
\item \texttt{coarse\_sand\_1} in \(\mathrm{g\,L^{-1}}\);
\item \texttt{mp1} in \(\mathrm{kg\,m^{-3}}\).
\end{itemize}

Two points are especially clear:
\begin{enumerate}
\item \texttt{mp1} is constant in time and identical for all 9 rivers:
\[
\texttt{mp1} = 1.754906975\times 10^{-6}\ \mathrm{kg\,m^{-3}}.
\]
\item \texttt{coarse\_sand\_1} varies in time, with domain-wide minimum and maximum values
\[
4.516\times 10^{-4}\ \mathrm{g\,L^{-1}} \le \texttt{coarse\_sand\_1} \le 1.6097\times 10^{-2}\ \mathrm{g\,L^{-1}}.
\]
\end{enumerate}

Because the plastic concentration is constant, the imposed plastic source differs among rivers only
through the discharge \texttt{river\_flux}. Using the hourly records in \texttt{Times}, the implied annual
river-source totals are approximately:
\begin{itemize}
\item Delaware group: \(1.379\times 10^{10}\ \mathrm{m^3}\) water,
\(2.420\times 10^{4}\ \mathrm{kg}\) plastic,
\(6.578\times 10^{7}\ \mathrm{kg}\) sediment;
\item Schuylkill group: \(3.915\times 10^{9}\ \mathrm{m^3}\) water,
\(6.870\times 10^{3}\ \mathrm{kg}\) plastic,
\(1.351\times 10^{7}\ \mathrm{kg}\) sediment.
\end{itemize}

\section{Case-by-Case Reconstruction}
\begin{longtable}{>{\raggedright\arraybackslash}p{0.08\textwidth} >{\raggedright\arraybackslash}p{0.24\textwidth} >{\raggedright\arraybackslash}p{0.18\textwidth} >{\raggedright\arraybackslash}p{0.18\textwidth} >{\raggedright\arraybackslash}p{0.18\textwidth}}
\toprule
Case & Intended role recovered from files & Plastic floc switch & Sediment setting & Other notable differences \\
\midrule
\endfirsthead
\toprule
Case & Intended role recovered from files & Plastic floc switch & Sediment setting & Other notable differences \\
\midrule
\endhead
\texttt{a} &
Baseline plastic transport case without active plastic--sediment floc interaction. &
\texttt{PLAST\_FLOC = F} in \texttt{generic\_plastic\_a.inp}. &
\texttt{SEDIMENT\_MODEL = F}. The sediment file name is still present in the namelist but should be inactive. &
Uses \texttt{STARTUP\_TS\_TYPE = 'observed'} from \texttt{waterPACT\_DRE\_its.nc}. NetCDF and restart output begin on MJD 58490.0, i.e. 5 days after the model start. \\
\texttt{b} &
Main floc-enabled Delaware test with active suspended sediment and river plastic input. &
\texttt{PLAST\_FLOC = T} in \texttt{generic\_plastic\_b.inp}. &
\texttt{SEDIMENT\_MODEL = T} with \texttt{generic\_sediment.inp}; the only sediment class is \texttt{coarse\_sand\_1} with \texttt{SED\_SD50 = 0.1 mm}. &
Now uses \texttt{STARTUP\_TS\_TYPE = 'observed'} from \texttt{waterPACT\_DRE\_its.nc}, matching case \texttt{a}. NetCDF and restart output begin at the model start. \\
\texttt{c} &
Sensitivity test relative to case \texttt{b}, most likely aimed at changing sediment character while keeping the plastic settings fixed. &
\texttt{PLAST\_FLOC = T}. \texttt{generic\_plastic\_c.inp} is byte-for-byte identical to \texttt{generic\_plastic\_b.inp}. &
\texttt{SEDIMENT\_MODEL = T} with \texttt{generic\_sediment\_c.inp}. The only actual change is \texttt{SED\_SD50 = 1.0 mm} instead of \(0.1\) mm. &
Same run configuration as case \texttt{b} except for output directory, sediment file name, plastic file name, and a shorter batch request on the cluster script. It now also uses \texttt{STARTUP\_TS\_TYPE = 'observed'}. \\
\bottomrule
\end{longtable}

\section{Important Non-Obvious Differences}
\subsection{The floc interaction term is explicitly sediment-size dependent}
The water-column floc interaction is built from two rate fields in \texttt{mod\_plast.F}:
\(\lambda_a\) for aggregation and \(\lambda_d\) for disaggregation. These are recomputed every step when
\texttt{PLAST\_FLOC = T}. The sediment median diameter \texttt{sed(ised)\%Sd50} enters the interaction in
several direct ways:
\begin{itemize}
\item In \texttt{Calc\_floc\_MP\_Aggregation\_Rate}, each sediment class defines a floc-size gate
\(d_{MP}^{max}\) from \texttt{Sd50}. If the plastic size exceeds this gate, that sediment class contributes
zero aggregation.
\item The same routine uses \texttt{Sd50} to define the sediment particle radius \(r_{sed}\), which enters
both the turbulent shear kernel \((r_{sed}+r_{plast})^3\) and the differential-settling kernel
\((r_{sed}+r_{plast})^2 |w_{sed}-w_{plast}|\).
\item The sediment mass concentration is converted to a number density using a spherical-particle mass
\(\propto r_{sed}^3\), so larger \texttt{Sd50} means fewer particles for the same sediment mass
concentration.
\item In \texttt{Calc\_floc\_MP\_Disaggregation\_Rate}, the maximum active sediment size across classes is used
to rebuild the same \(d_{MP}^{max}\) gate, which then appears in the disaggregation coefficient.
\end{itemize}

\subsection{Current consequence for the Delaware test cases}
For the present Delaware inputs, the plastic particle size is approximately \(0.100\ \mathrm{mm}\)
\(=100\ \mu\mathrm{m}\). The sediment-size gate therefore behaves very differently in cases
\texttt{b} and \texttt{c}:
\begin{itemize}
\item Case \texttt{b}: \texttt{SED\_SD50 = 0.1 mm = 100 \(\mu\)m}. The implemented gate gives
\[
d_{MP}^{max} = -0.0002\,d_{floc}^2 + 0.36\,d_{floc} \approx 34\ \mu\mathrm{m},
\]
so a \(100\ \mu\mathrm{m}\) microplastic is larger than the permitted size. In practice this means the
aggregation kernel is strongly size-excluded.
\item Case \texttt{c}: \texttt{SED\_SD50 = 1.0 mm = 1000 \(\mu\)m}. The code then uses the large-floc cap
\[
d_{MP}^{max} = 162\ \mu\mathrm{m},
\]
so the same \(100\ \mu\mathrm{m}\) microplastic is now allowed to aggregate.
\end{itemize}

This is not a small secondary effect. It means case \texttt{b} may behave close to a no-aggregation limit,
while case \texttt{c} is the first case in this bundle that can genuinely exercise the floc aggregation
kinetics for the chosen \(100\ \mu\mathrm{m}\) plastic.

\subsection{Aggregation changes the settling branch as well as the split state}
The same code path assigns different settling velocities to the two plastic sub-states:
\begin{itemize}
\item \texttt{conc\_d} uses \texttt{Calc\_Wset\_Plastic}, which for the present
\(\rho_p = 950\ \mathrm{kg\,m^{-3}}\) plastic in seawater gives a negative settling velocity
(upward or buoyant motion);
\item \texttt{conc\_a} is settled with the sediment settling velocity from the sediment model
(\texttt{sed(1)\%Wset} in the current implementation).
\end{itemize}

Therefore the floc interaction does not merely repartition \texttt{conc\_a} and \texttt{conc\_d}; once
aggregation is active it can switch plastic from a buoyant branch to a downward-settling branch.

\subsection{The startup T/S mismatch has now been removed}
The run namelists for cases \texttt{b} and \texttt{c} have now been updated so that all three cases use
\texttt{STARTUP\_TS\_TYPE = 'observed'} with the same startup file \texttt{waterPACT\_DRE\_its.nc}. The
startup file contains nontrivial vertical structure:
\begin{itemize}
\item salinity \texttt{ssl} ranges from about \(-0.0028\) to \(35.19\) PSU;
\item temperature \texttt{tsl} ranges from about \(4.82\) to \(13.43^\circ\mathrm{C}\).
\end{itemize}

This removes one physical confounder from the original comparison between cases \texttt{a}, \texttt{b},
and \texttt{c}. A run-control difference still remains in the output timing:
\begin{itemize}
\item \texttt{RST\_FIRST\_OUT}: day 58490.0 in \texttt{a}, day 58485.0 in \texttt{b}/\texttt{c};
\item \texttt{NC\_FIRST\_OUT}: day 58490.0 in \texttt{a}, day 58485.0 in \texttt{b}/\texttt{c}.
\end{itemize}

\subsection{Case \texttt{c} only changes sediment grain size, not the plastic setup}
The \texttt{generic\_plastic\_b.inp} and \texttt{generic\_plastic\_c.inp} files are identical. The only
content difference recovered between \texttt{b} and \texttt{c} is:
\[
\texttt{SED\_SD50}: 0.1\ \mathrm{mm} \rightarrow 1.0\ \mathrm{mm}.
\]

Other sediment parameters such as \texttt{SED\_WSET}, \texttt{SED\_TAUE}, and \texttt{SED\_TAUD} remain
unchanged in the provided files. So case \texttt{c} is best interpreted as a grain-size sensitivity test,
not a broader sediment physics retuning.

\subsection{No preserved model outputs are available in the case output folders}
The directories \texttt{OUTPUT\_a}, \texttt{OUTPUT\_b}, and \texttt{OUTPUT\_c} are presently empty. That means
the original scenario intent has to be inferred entirely from the inputs and batch scripts; there is no
saved output available here to confirm what diagnostics were originally examined.

\section{Case \texttt{b} Output NetCDF Structure}
The first successful case-\texttt{b} output stack inspected on Kestrel is:
\[
\texttt{/kfs3/scratch/yhuang168/waterPACT\_MP\_floc/OUTPUT\_b/waterPACT\_b\_0001.nc}.
\]
It is a 64-bit NetCDF file with 160 records on the unlimited \texttt{time} dimension. A fuller
script-oriented inventory has been saved separately as
\texttt{analysis\_notes/waterPACT\_netcdf\_structure\_b.md}. The no-floc baseline structure for
experiment \texttt{a} is recorded in
\texttt{analysis\_notes/waterPACT\_netcdf\_structure\_a.md}.

\subsection{Core Dimensions}
\begin{center}
\begin{tabular}{lr}
\toprule
Dimension & Size \\
\midrule
\texttt{nele} & 126722 \\
\texttt{node} & 68416 \\
\texttt{siglay} & 30 \\
\texttt{siglev} & 31 \\
\texttt{three} & 3 \\
\texttt{time} & 160, unlimited \\
\texttt{maxnode} & 11 \\
\texttt{maxelem} & 9 \\
\texttt{four} & 4 \\
\bottomrule
\end{tabular}
\end{center}

\subsection{Diagnostic-Relevant Variables}
The fields most relevant for later mass-budget scripts are:
\begin{itemize}
\item geometry and volumes: \texttt{art1(node)}, \texttt{h(node)}, \texttt{zeta(node,time)},
\texttt{siglay(node,siglay)}, and \texttt{siglev(node,siglev)};
\item sediment forcing/state: \texttt{coarse\_sand\_1(node,siglay,time)} and
\texttt{coarse\_sand\_1\_bedfrac(node,time)};
\item total suspended plastic: \texttt{mp1(node,siglay,time)};
\item floc split: \texttt{mp1\_agg(node,siglay,time)} and
\texttt{mp1\_dis(node,siglay,time)};
\item interaction rates and settling: \texttt{lambda\_a\_mp1}, \texttt{lambda\_d\_mp1},
\texttt{settle\_vel\_mp1}, and \texttt{settle\_vel\_floc\_mp1};
\item bed exchange and bed inventory: \texttt{dep\_flux\_mp1(node,time)},
\texttt{ero\_flux\_mp1(node,time)}, \texttt{bot\_mass\_mp1(node,time)},
\texttt{bot\_mass\_agg\_mp1(node,time)}, and \texttt{bot\_mass\_dis\_mp1(node,time)};
\item supporting stress/density fields: \texttt{taub\_total(node,time)}, \texttt{bottom\_stress(node,time)},
\texttt{tau\_ce\_mp1(node,time)}, and \texttt{ambient\_rho(node,siglay,time)}.
\end{itemize}

For script design, the main consistency check is:
\[
\texttt{mp1} \approx \texttt{mp1\_agg} + \texttt{mp1\_dis}.
\]
Water-column inventory should use node-control-volume areas, likely \texttt{art1}, together with
time-varying water depth and sigma-layer thickness. Bed inventory should use area integration of
\texttt{bot\_mass\_*} fields. For the no-floc baseline case \texttt{a}, floc-only output variables may be
absent after the NetCDF-registration fix because \texttt{PLAST\_FLOC = F}.

\section{Expected Mass-Conserving Signatures}
If the model is internally mass conservative, the comparison among \texttt{a}, \texttt{b}, and \texttt{c}
should look like a redistribution problem rather than a creation/loss problem.

\subsection{Closed-system expectation}
In a closed or internally isolated test, the correct quantity to compare is the total plastic inventory
\[
M_{tot} = M_{water} + M_{bed}.
\]
For all three cases, \(M_{tot}\) should remain constant up to numerical truncation. The floc interaction
should only move plastic among:
\begin{enumerate}
\item disaggregated water-column plastic \texttt{conc\_d},
\item aggregated water-column plastic \texttt{conc\_a},
\item bottom or bed plastic mass.
\end{enumerate}

\subsection{What case \texttt{a} versus case \texttt{b} should look like}
Because the current size gate in case \texttt{b} appears to exclude aggregation for a
\(100\ \mu\mathrm{m}\) microplastic in \(100\ \mu\mathrm{m}\) sediment, case \texttt{b} should be close to
the no-floc baseline in its floc kinetics:
\begin{itemize}
\item \texttt{conc\_a} should remain near zero or at least much smaller than \texttt{conc\_d};
\item the total plastic inventory in case \texttt{b} should closely track case \texttt{a} when summed over
water plus bed;
\item if case \texttt{b} shows a large total-mass deficit relative to case \texttt{a}, that would be more
consistent with a numerical leak than with the intended floc physics.
\end{itemize}

\subsection{What case \texttt{a} versus case \texttt{c} should look like}
Case \texttt{c} is the one that should reveal the sediment-size sensitivity most clearly:
\begin{itemize}
\item \texttt{conc\_a} should become nonzero where suspended sediment is present;
\item \texttt{conc\_d} should decrease where aggregation is active;
\item the water-column plastic inventory can drop relative to case \texttt{a}, but that drop should be
matched by an increase in bed plastic mass rather than by a loss of total plastic;
\item because aggregated plastic uses the sediment settling branch while disaggregated plastic is buoyant,
case \texttt{c} should show a stronger downward redistribution than case \texttt{a}.
\end{itemize}

\subsection{What case \texttt{b} versus case \texttt{c} should look like}
If the floc model is working and mass is conserved:
\begin{itemize}
\item case \texttt{c} should have more \texttt{conc\_a} and more bed plastic than case \texttt{b};
\item case \texttt{c} should have less floating or disaggregated plastic than case \texttt{b};
\item the difference should be mainly in partitioning and vertical redistribution, not in total
\(M_{tot}\).
\end{itemize}

\subsection{Open-boundary caveat}
With realistic river and open-boundary forcing active, domain-integrated water-column mass does not need to
match among cases at every output time because aggregation can change residence time and export. Even so, an
internally conservative solution should still satisfy:
\begin{itemize}
\item no unexplained jump in \(\texttt{conc\_a}+\texttt{conc\_d}\) across the kinetics step itself;
\item no persistent loss of \(M_{water}+M_{bed}\) that cannot be explained by diagnosed export through open
boundaries.
\end{itemize}

\section{Supporting Scripts}
Reusable inspection utilities have been placed in \texttt{FVCOM\_MP\_test\_run/PYTHON}. The main script
\texttt{inspect\_delaware\_cases.py} writes:
\begin{itemize}
\item \texttt{FVCOM\_MP\_test\_run/PYTHON/output/case\_inventory.json};
\item \texttt{FVCOM\_MP\_test\_run/PYTHON/output/case\_inventory.md}.
\end{itemize}

These files provide a machine-readable and human-readable summary of the same recovered scenario information.

\section{Bottom Line}
The recovered intent is most consistent with:
\begin{itemize}
\item \texttt{a}: no-floc baseline;
\item \texttt{b}: floc-enabled reference case;
\item \texttt{c}: floc-enabled sediment-size sensitivity case.
\end{itemize}

The most important code-level finding is that the floc interaction depends directly on sediment size, and
with the present inputs that dependence likely suppresses aggregation in case \texttt{b} while allowing it
in case \texttt{c}. After aligning startup T/S across all three run namelists, the main physical comparison
to watch is whether case \texttt{c} shows redistribution into \texttt{conc\_a} and the bed without losing
total plastic mass.

\end{document}
